\title{FlashAttention-3: Fast and Accurate Attention with Asynchrony and Low-precision}

\begin{abstract}
Attention mechanisms, integral to the foundational Transformer architecture, represent a primary computational constraint in scaling large language models and managing extended input sequences. \textsc{FlashAttention} introduced a strategy for accelerating attention on GPUs by significantly reducing memory access. Nevertheless, this method does not fully leverage advancements in modern GPU hardware; for instance, \textsc{FlashAttention-2} only utilizes 35\% of the H100 GPU's capacity. To further optimize performance on Hopper GPUs, we introduce three key enhancements: (1) harnessing the asynchronous capabilities of Tensor Cores and TMA to enable warp-specialized parallelism, thereby allowing computation and data transfer to proceed concurrently; (2) integrating block-wise matrix multiplication and softmax computation for improved interleaving; and (3) employing block quantization and disjoint processing through hardware-enabled FP8 support for low-precision arithmetic. Our proposed solution, \textsc{FlashAttention-3}, demonstrates 1.5-2.0$\times$ speedup on H100 GPUs, attaining up to 740 TFLOPs/s with FP16 (reflecting 75\% utilization), and nearing 1.2 PFLOPs/s in FP8 precision. Moreover, we show that FP8 \textsc{FlashAttention-3} offers a 2.6$\times$ reduction in numerical error compared to conventional FP8 attention baselines.
\end{abstract}

\section{Introduction}
\label{sec:intro}

Within the Transformer architecture~\citep{vaswani2017attention}, the principal computational constraint arises from the attention mechanism, primarily because calculating self-attention scores between queries and keys results in a quadratic time complexity with respect to the input sequence length. Enhancing attention mechanisms to support much longer contexts can facilitate novel functionalities, such as handling and reasoning over several extended documents~\citep{guo2021longt5,shaham2022scrolls,peng2023yarn} or navigating extensive code repositories~\citep{roziere2023code, li2023starcoder}. Additionally, it paves the way for processing new data types, including high-resolution images~\citep{chen2022scaling}, audio streams~\citep{gulati2020conformer}, and video sequences~\citep{ho2022video}, as well as powering applications requiring long-term user interaction~\citep{sun2019bert4rec} or agent tasks with extended horizons~\citep{yao2022react}. Consequently, there is considerable ongoing research aimed at accelerating attention computations in long-context settings through methods such as attention approximations~\citep{katharopoulos2020transformers,choromanski2020rethinking,
tay2020efficient}, improved software implementations (\citep{rabe2021self,dao2022flashattention,kwon2023efficient}), and even by investigating novel model architectures~\citep{peng2023rwkv,sun2023retentive,gu2023mamba}.

This study extends the advancements of \citet{dao2022flashattention}, who formulated precise attention computation methods by embedding insights about GPU architecture and its underlying operational patterns into their algorithmic frameworks. Dao et al. in \citep{dao2022flashattention} introduced \textsc{FlashAttention}, a novel approach utilizing a tiling mechanism to parallelize attention computation; their method consolidated all steps of the attention process into a singular GPU kernel, thereby avoiding the need for slow global memory accesses due to intermediate data movement. 

Subsequently, \citet{dao2023flashattention2} revised the method as \textsc{FlashAttention-2}, incorporating parallelization along the sequence length axis and executing the inner forward-pass loop across segmented blocks of the key and value matrices. This enhancement aimed to boost both workload balance and resource utilization on the GPU. Nonetheless, we find that \textsc{FlashAttention-2} demonstrates significantly lower resource usage on state-of-the-art GPUs compared to high-performance matrix multiplication (GEMM) kernels; specifically, the method attains approximately 35\% utilization, whereas GEMM can reach between 80–90\% on the Hopper H100 GPU. This performance gap may be partly due to implementation-specific factors, including reliance on Ampere instructions rather than those tailored for the Hopper architecture when running on Tensor Cores. Recent work such as ThunkerKitten~\citep{spector2024thunder} and cuDNN 9~\citep{cudnn9} indicates that leveraging Hopper-optimized instructions in conjunction with tile-based computational paradigms can accelerate attention operations and reduce implementation complexity.

At a deeper level, the algorithm behind \textsc{FlashAttention-2} operates within a basic synchronous paradigm and does not explicitly incorporate asynchronous execution or low-precision arithmetic in its architecture. Asynchronous execution emerges due to hardware components optimized for key machine learning tasks: dedicated units such as Tensor Cores execute matrix multiplications, while the Tensor Memory Accelerator (TMA) handles memory operations, functioning independently from the main CUDA cores responsible for general computation involving logic, integer, and floating point operations. Progress in low-precision computing—manifested in FP8 formats on Hopper and FP4 on Blackwell, following earlier developments like FP16 on Pascal (2017) and BF16 on Ampere (2020)—has consistently allowed for significant improvements in throughput per unit area and power. The relevant hardware features provided by Hopper along these dimensions are discussed in \cref{subsec:hardware}.  
The principal technical task lies in adapting \textsc{FlashAttention-2} so that it leverages these hardware capabilities: supporting asynchrony necessitates orchestrating overlapping execution of matmul and softmax even with their inherent data dependencies, while adopting low-precision computations demands careful strategies to mitigate quantization artifacts, particularly for extreme-value features present in LLMs~\citep{dettmers2208llm, sun2024massive}.

With this objective, we introduce \textsc{FlashAttention-3}, integrating and building upon three novel concepts aimed at enhancing efficiency on modern GPU architectures.\footnote{Although our findings are contextualized within NVIDIA's Hopper architecture, the proposed algorithm is applicable to any GPU platform that supports strong asynchronous execution and low-precision computation.}

\begin{enumerate}

\item \textbf{Producer-Consumer asynchrony:} We introduce a warp-specialized software pipelining approach that capitalizes on the ability to execute data transfers and Tensor Core operations independently. By allocating producer and consumer responsibilities to distinct warps, our scheme extends the capability of the algorithm to mask both memory and instruction issuance latencies.
\item \textbf{Hiding softmax under asynchronous block-wise GEMMs:} To optimize overlap, we interleave the execution of relatively low-throughput softmax operations—such as floating-point fused multiply-adds and exponentials—with asynchronous GEMM computations via WGMMA instructions. This redesign involves modifying the \textsc{FlashAttention-2} algorithm to relax sequential dependencies between softmax computation and GEMM operations. In our two-stage variant, for instance, the softmax function is computed on one section of the score matrix while, simultaneously, WGMMA asynchronously processes the subsequent section.
\item \textbf{Hardware-accelerated low-precision GEMM:} Our forward pass implementation is revised to leverage FP8 Tensor Cores during GEMM operations, achieving close to twice the observed TFLOPs/s throughput. Addressing the disparity in layout assumptions between WGMMA (for FP32 accumulators and FP8 input matrices) and standard storage patterns, we employ block quantization and incoherent processing strategies. These techniques help compensate for potential accuracy degradation stemming from the use of FP8 arithmetic.
\end{enumerate}

To empirically assess our approach, we conduct benchmarks of \textsc{FlashAttention-3} on an H100 SXM5 GPU, spanning various parameter configurations. Our results indicate that (1) in FP16, the forward pass exhibits a 1.5–2.0$\times$ acceleration over \textsc{FlashAttention-2} (peaking at 740 TFLOPs/s), while the backward pass attains a 1.5–1.75$\times$ speedup; (2) FP8 performance approaches 1.2 PFLOPs/s; and (3) at longer sequence lengths, FP16 outperforms alternatives, and FP8 demonstrates competitiveness\footnote{Specifically, with a head dimension of 64, FP8 \textsc{FlashAttention-3} surpasses others, whereas for head dimensions 128 and 256, it matches alternatives in non-causal settings but lags behind when causal masking is used.} with NVIDIA’s cuDNN library’s top-performing attention implementation. Further evaluation confirms that FP16 \textsc{FlashAttention-3} matches \textsc{FlashAttention-2} in numerical precision, and surpasses traditional attention—attributable to retaining intermediates, such as the rescaled softmax, in FP32. Additionally, FP8 \textsc{FlashAttention-3}, leveraging block quantization and incoherent computation, achieves 2.6$\times$ higher accuracy compared to conventional attention with per-tensor quantization, particularly in the presence of outlier features.

\textsc{FlashAttention-3} has been released under an open and permissive license\footnote{\textsc{FlashAttention-3}
  is available at \texttt{https://github.com/Dao-AILab/flash-attention}}. We also aim to incorporate it into the PyTorch and Hugging Face ecosystems to maximize its accessibility for researchers and practitioners.

\section{Background: Multi-Head Attention and GPU Characteristics}
\label{sec:background}

\subsection{Multi-Head Attention}
\label{subsec:multi_head_attn}

Consider $\mathbf{Q}, \mathbf{K}, \mathbf{V} \in \mathbb{R}^{N \times d}$, representing the query, key, and value matrices for a particular attention head, where $N$ denotes the length of the input sequence and $d$ corresponds to the dimensionality of each head. The resulting attention output, $\mathbf{O}$, is then given by:

\begin{equation*}
  \mathbf{S} = \alpha \mathbf{Q} \mathbf{K}^\top \in \mathbb{R}^{N \times N}, \quad \mathbf{P} = \mathrm{softmax}(\mathbf{S}) \in \mathbb{R}^{N \times N}, \quad \mathbf{O} = \mathbf{P}\mathbf{V} \in \mathbb{R}^{N \times d},
\end{equation*}

Here, $\mathrm{softmax}$ operates along each row, and it is common to choose $\alpha = 1/\sqrt{d}$ as the scaling parameter. To address potential numerical issues due to the exponential, it is standard practice to subtract $\mathrm{rowmax}(\mathbf{S})$ from $\mathbf{S}$. In the multi-head attention (MHA) setting, separate projections for queries, keys, and values are used for each head. This allows the computations to be performed in parallel across heads and batches, resulting in the aggregate output tensor.

Consider a scalar-valued loss function $\phi$, and denote the gradient with respect to an argument by $\mathbf{d}(-) = \partial \phi / \partial (-)$. Provided the gradient with respect to the output, $\mathbf{dO} \in \mathbb{R}^{N \times d}$, the gradients for $\mathbf{Q}$, $\mathbf{K}$, and $\mathbf{V}$—namely $\mathbf{dQ}$, $\mathbf{dK}$, and $\mathbf{dV}$—are derived via the chain rule as shown below:

\begin{align*}
  \mathbf{dV} &= \mathbf{P}^\top \mathbf{dO} \in \mathbb{R}^{N \times d} \\
  \mathbf{dP} &= \mathbf{dO} \mathbf{V}^\top \in \mathbb{R}^{N \times N} \\
  \mathbf{dS} &= \mathrm{dsoftmax} (\mathbf{dP}) \in \mathbb{R}^{N \times N} \\
  \mathbf{dQ} &= \alpha \mathbf{dS} \mathbf{K} \in \mathbb{R}^{N \times d} \\
  \mathbf{dK} &= \alpha \mathbf{dS}^\top \mathbf{Q} \in \mathbb{R}^{N \times d},
\end{align*}

In this context, $\mathbf{d}s = (\mathrm{diag}(p) - p p^\top)\mathbf{d}p$ holds, where $p = \mathrm{softmax}(s)$ is treated as a function dependent on the vector $s$. The notation $\mathrm{dsoftmax}(\mathbf{dP})$ indicates the application of this expression on a row-by-row basis. Moreover, this operation can be efficiently parallelized over both the number of heads and the batch dimension during the backward propagation step in MHA.

\subsection{GPU hardware characteristics and execution model}
\label{subsec:hardware}

We discuss elements of the GPU execution model pertinent to \textsc{FlashAttention-3}, emphasizing the NVIDIA Hopper architecture as a specific example of this framework.

\paragraph{Memory hierarchy:} The memory system of the GPU is structured as a hierarchy of storage locations, where higher-capacity components offer lower bandwidth, as summarized in \cref{tab:gpu-hierarchy}\footnote{\citet{luo2024benchmarking} states that shared memory achieves 128 bytes per clock cycle per SM; by scaling with 132 SMs and the 1830 MHz boost clock, overall bandwidth is estimated.}. At the base, global memory (GMEM), or HBM, serves as the off-chip DRAM shared by all streaming multiprocessors (SMs). Data stored in GMEM is automatically cached in the on-chip L2 cache to improve access speed. Each SM is also equipped with a relatively small, highly banked, on-chip shared memory (SMEM) that is directly managed by the programmer. At the top of the hierarchy lies the register file, found within every SM.

\paragraph{Thread hierarchy:} In the GPU programming paradigm, execution units are systematically arranged into a hierarchical structure referred to as threads. The organization spans several levels, beginning with individual threads at the lowest granularity, followed by warps (each consisting of 32 threads), warpgroups (each made up of 4 consecutive warps), and progressing to threadblocks (also known as cooperative thread arrays or CTAs), threadblock clusters (introduced in Hopper architectures), and finally, grids at the highest level.

The two hierarchies exhibit a strong interconnection. Threads belonging to a single CTA are concurrently executed on a shared SM, while CTAs grouped in the same cluster are scheduled together on an identical GPC. All threads within a CTA have direct access to SMEM, whereas individual threads possess up to 256 private registers (RMEM) exclusively for their own use.

\paragraph{Asynchrony and warp-specialization:}

GPUs function as throughput-oriented processors, leveraging parallelism and asynchronous operations to mask the delays associated with memory accesses and execution. In the Hopper architecture, a specialized hardware component known as the Tensor Memory Accelerator (TMA) handles asynchronous memory transfers between GMEM and SMEM \cite[\S7.29]{cuda}. Additionally, in contrast to previous designs like Ampere, Hopper's Tensor Core—accessed through the warpgroup-level WGMMA instruction \cite[\S9.7.14]{ptx}—operates asynchronously and is capable of directly retrieving its input data from shared memory.

Asynchronous hardware capabilities facilitate the implementation of warp-specialized kernels, wherein warps within a CTA are partitioned into producer and consumer groups dedicated solely to data movement or computation tasks, respectively. This division typically enhances the compiler’s capacity to produce efficient instruction schedules \citep{warp-specialization-2011}. Moreover, Hopper introduces the ability to dynamically redistribute registers among warpgroups using \verb|setmaxnreg| \cite[\S9.7.17.1]{ptx}. This permits warps engaged in MMAs to access a greater portion of RMEM compared to those handling TMA operations, which require only a single thread.

\paragraph{Low-precision number formats:}
\label{sec:low-precision-gpu}
Contemporary GPUs are equipped with dedicated hardware modules that facilitate efficient low-precision arithmetic. For instance, on the Hopper architecture, employing the WGMMA instruction leverages the FP8 Tensor Cores, resulting in twice the throughput per SM relative to computations using FP16 or BF16 formats.

Successfully utilizing FP8 WGMMA requires a clear understanding of the operand layout requirements. When executing a GEMM operation computing $A \times B^{\top}$, with $A$ of dimensions $M\times K$ and $B$ of dimensions $N\times K$, the operand $A$ or $B$ is described as \emph{mn-major} if its memory layout is contiguous along the outer $M$ or $N$ dimension, and as \emph{k-major} if contiguity is along the inner $K$ dimension. In the context of FP16 WGMMA, either mn-major or k-major data arrangements are permissible for inputs stored in SMEM. In contrast, FP8 WGMMA restricts input operands to the k-major organization. This restriction becomes problematic in complex scenarios—for example, attention mechanisms—where consecutive GEMM operations are fused within a single kernel. Here, the mismatched layouts between FP32 accumulators and FP8 inputs hinder the chaining of FP8 WGMMA calls.

With regard to attention mechanisms, these layout constraints necessitate specific adaptations to the FP8 algorithm's architecture, as outlined in \cref{sec:algofp8}.

\subsection{Standard Attention and Flash Attention} In line with \citet{dao2022flashattention}, we use the term \textbf{standard attention} to refer to a GPU-based attention implementation that stores the intermediate matrices $\mathbf{S}$ and $\mathbf{P}$ in high-bandwidth memory (HBM). \textsc{FlashAttention} introduces a key optimization: by applying a localized softmax reduction, it circumvents the need for costly intermediate memory operations and consolidates the attention computation into a unified kernel. This local softmax process is realized in lines \ref{code-ws:softmax_start}-\ref{code-ws:softmax_end} of the consumer mainloop in \cref{alg:flash3_wgmma_ws_only}, as well as through the rescaling of blocks within $\mathbf{O}$. For a concise proof that this strategy accurately computes $\mathbf{O}$, see \cite[\S 2.3.1]{dao2023flashattention2}.

\section{FlashAttention-3: Algorithm}
\label{sec:algo}

This section outlines the \textsc{FlashAttention-3} algorithm. Our exposition centers on the forward computation; details of the backward algorithm can be found in~\cref{sec:algo_ws_bwd}. We begin by showing how warp-specialization, paired with a circular SMEM buffer, can be incorporated into the core procedure of \textsc{FlashAttention-2}. Next, we discuss leveraging the asynchronous nature of WGMMA to construct a two-stage pipeline that overlaps GEMM and softmax computations. Lastly, we cover the adjustments required to support FP8, addressing both layout compliance and accuracy enhancements through block quantization and incoherent processing.

\subsection{Producer-Consumer asynchrony through warp-specialization and pingpong scheduling}
\label{sec:algo_ws}

\paragraph{Warp-specialization}
Similar to \textsc{FlashAttention-2}, the forward computation in \textsc{FlashAttention-3} can be performed in a highly parallel manner across the batch dimension, head indices, and the length of the query sequence. Therefore, it is sufficient to describe the algorithm from the perspective of a CTA, focusing on its operation on a specific query matrix tile $\mathbf{Q}_i$ in order to produce the corresponding output tile $\mathbf{O}_i$. For clarity, we initially describe the warp-specialization approach utilizing a circular SMEM buffer, intentionally omitting the GEMM-softmax overlap mechanism. Let $d$ denote the head dimension and $N$ the length of the sequence; select a query block size $B_r$ such that $\mathbf{Q}$ is partitioned into $T_r = \lceil \frac{N}{B_r} \rceil$ segments, labeled as $\mathbf{Q}_1, .., \mathbf{Q}_{T_r}$.

\begin{algorithm}[H]
    \caption{\small\label{alg:flash3_wgmma_ws_only}\textsc{FlashAttention-3} forward pass \textbf{excluding} intra-consumer overlap -- CTA perspective}
    \begin{algorithmic}[1]
\REQUIRE Matrices $\mathbf{Q}_i \in \mathbb{R}^{B_r \times d}$, and $\mathbf{K}, \mathbf{V} \in \mathbb{R}^{N \times d}$ reside in HBM; use key block size $B_c$, define $T_c = \lceil \frac{N}{B_c} \rceil$.
\STATE Set up a pipeline handler to control barrier synchronization using an $s$-stage circular buffer in shared memory (SMEM).
\IF {assigned to the producer warpgroup}
\STATE Free a pre-specified number of registers.
\STATE Transfer $\mathbf{Q}_i$ from HBM to shared memory.
\STATE After loading completes, notify consumers of the availability of $\mathbf{Q}_i$.
\FOR{$0 \le j < T_c$}
    \STATE Ensure the $(j\,\%\,s)$-th buffer stage has been consumed.
    \STATE Load $\mathbf{K}_j, \mathbf{V}_j$ from HBM into the $(j\,\%\,s)$ stage of the shared memory buffer.
    \STATE On completion, signal to consumers that $\mathbf{K}_j, \mathbf{V}_j$ are ready.
\ENDFOR
\ELSE \STATE Reassign a pre-established quantity of registers, determined by the count of consumer warps.
\STATE On-chip, initialize outputs $\mathbf{O}_i = (0) \in \mathbb{R}^{B_r \times d}$ and the accumulator variables $\ell_i, m_i = (0), (-\infty) \in \mathbb{R}^{B_r}$.
\STATE Wait for $\mathbf{Q}_i$ to finish loading into shared memory.
\FOR{$0 \le j < T_c$}
\STATE Wait until $\mathbf{K}_j$ becomes available in shared memory.
\STATE Evaluate $\mathbf{S}_i^{(j)} = \mathbf{Q}_i \mathbf{K}_j^T$ via SS-GEMM. Commit and synchronize as needed. \label{alg:ws_only_gemm1}
\STATE Set $m_i^{\mathrm{old}} = m_i$, then update $m_i = \mathrm{max}(m_i^{\mathrm{old}}, \mathrm{rowmax}(\mathbf{S}_i^{(j)}))$. \label{code-ws:softmax_start}
\STATE Calculate $\widetilde{\mathbf{P}}_i^{(j)} = \mathrm{exp}(\mathbf{S}_i^{(j)} - m_i)$ and update $\ell_i = \mathrm{exp}(m_i^{\mathrm{old}} - m_i) \ell_i + \mathrm{rowsum}(\widetilde{\mathbf{P}}_i^{(j)})$. \label{code-ws:softmax_end}
\STATE Wait for the $\mathbf{V}_j$ data to be ready in shared memory.
\STATE Compute $\mathbf{O}_i = \mathrm{diag}(\mathrm{exp}(m_i^{\mathrm{old}} - m_i))^{-1} \mathbf{O}_i + \widetilde{\mathbf{P}}_i^{(j)} \mathbf{V}_j$ using RS-GEMM. Commit and synchronize as needed. \label{alg:ws_only_gemm2}
\STATE Mark the $(j\,\%\,s)$th stage as available for producer use.
\ENDFOR
\STATE Finalize $\mathbf{O}_i = \mathrm{diag}(\ell_i)^{-1} \mathbf{O}_i$ and compute $L_i = m_i + \log(\ell_i)$.
\STATE Save the resulting $\mathbf{O}_i$ and $L_i$ to HBM, storing them as the $i$th segments of $\mathbf{O}$ and $L$.
\ENDIF
\end{algorithmic}
\end{algorithm}

In our realization of \cref{alg:flash3_wgmma_ws_only} on Hopper, we employ \verb|setmaxnreg| to manage (de)allocations, utilize TMA to fetch $\mathbf{Q}_i$ as well as $\{ \mathbf{K}_j, \mathbf{V}_j \}_{0 \leq j < T_c}$, and leverage WGMMA for carrying out the GEMM operations within the consumer mainloop. Here, either the SS or RS prefix specifies whether the initial operand is obtained from shared memory or the register file, respectively. It is important to recognize, when analyzing the execution order of \cref{alg:flash3_wgmma_ws_only}, that TMA load instructions are non-blocking—meaning their dispatch does not depend on the completion of other loads because of inherent asynchronicity. Additionally, throughout the producer mainloop, no explicit waits occur during the first $s$ iterations since the buffer is initially being populated.

\paragraph{Pingpong scheduling} The asynchrony enabled by WGMMA and TMA, in conjunction with warp specialization, allows for overlapping the softmax calculation performed by one warpgroup with the GEMM executed by another. This scheduling approach is motivated by the significant disparity in throughput between matmul and non-matmul operations on current hardware accelerators. For instance, the H100 SXM5 GPU is capable of delivering 989 TFLOPS in FP16 matrix multiplication, but achieves only 3.9 TFLOPS when executing special function operations such as the exponential\footnote{According to the CUDA programming guide, a streaming multiprocessor (SM) can process 16 special function operations per clock cycle. Multiplying 16 by 132 SMs and an 1830 MHz clock frequency yields 3.9 TFLOPS for special functions.}. In the context of attention forward passes with FP16 and a head dimension of 128, matrix multiplications require 512 times more FLOPS compared to exponentials, yet the throughput for exponentials is 256 times lower, implying that the exponential operation can consume as much as 50\% of the execution cycles required for matmul. This imbalance is exacerbated for FP8 computations, as matmul throughput doubles while exponential throughput does not increase.

Because the exponential operation is handled by a dedicated hardware component—the multi-function unit—it is most efficient to overlap its computation with the matrix multiplication performed by the Tensor Cores. To coordinate this, we employ synchronization barriers (\texttt{bar.sync} instructions), which ensure that the GEMM operations (specifically, GEMM1—$\mathbf{P} \mathbf{V}$ in a given iteration, and GEMM0—$\mathbf{Q} \mathbf{K}^\top$ in the subsequent iteration) in warpgroup 1 are completed before those in warpgroup 2 commence. This scheduling allows the softmax computation for warpgroup 1 to proceed concurrently with the GEMM executions in warpgroup 2. The process is then reversed: warpgroup 2 computes the softmax while warpgroup 1 returns to GEMMs, creating a “pingpong” execution pattern as shown in~\cref{fig:pingpong_scheduling}. While the real-world implementation of pingpong scheduling deviates somewhat from this schematic ideal, empirical results show noticeable performance benefits (for instance, FP16 forward throughput rises from 570 TFLOPS to between 620 and 640 TFLOPS with a head dimension of 128 and sequence length of 8192).

\paragraph{Attention variants} In the cases of multi-query attention~\citep{shazeer2019fast} and grouped query attention~\citep{ainslie2023gqa}, we adopt the strategy from \textsc{FlashAttention-2}, modifying tensor indices so as to prevent redundant storage of $\mathbf{K}$ and $\mathbf{V}$ in HBM.

\subsection{Intra-warpgroup overlapping GEMMs and softmax}

It is possible to achieve partial instruction overlap between softmax and GEMM operations within a single warpgroup. Here, we outline a specific method for accomplishing this.

Within the attention mechanism, certain steps in the core loop are inherently sequential, which restricts the extent of parallelism achievable in a single pass. For instance, the (local) softmax computation (lines \ref{code-ws:softmax_start} to \ref{code-ws:softmax_end}) is contingent on receiving $\mathbf{S}_i^{(j)}$ from the initial GEMM, and the subsequent GEMM utilizes the intermediate output $\widetilde{\mathbf{P}}_i^{(j)}$ as input. The explicit synchronization enforced by the wait statements at lines \ref{alg:ws_only_gemm1} and \ref{alg:ws_only_gemm2} in \cref{alg:flash3_wgmma_ws_only} effectively serializes the softmax and GEMM computations. Nonetheless, such dependencies may be decoupled by introducing pipeline parallelism between loop iterations, utilizing additional register buffers. Building on this insight, we present a pipelined GEMM-softmax strategy based on a two-stage design,\footnote{It is important to note that while the pipeline’s overlap stages are limited by, they do not necessarily match, the circular SMEM buffer’s $s$ stages.} as follows:

\begin{algorithm}[H]
    \caption{\small\label{alg:flash3_wgmma}\textsc{FlashAttention-3} consumer warpgroup forward pass}
    \begin{algorithmic}[1]
      \REQUIRE  Given matrices $\mathbf{Q}_i \in \mathbb{R}^{B_r \times d}$ as well as $\mathbf{K}, \mathbf{V} \in \mathbb{R}^{N \times d}$ located in HBM, and a specified key block size $B_c$, define $T_c = \lceil \frac{N}{B_c} \rceil$.
      \STATE Dynamically reassign a fixed quantity of registers according to the number of consumer warps engaged.
      \STATE On-chip initialization of $\mathbf{O}_i = (0) \in \mathbb{R}^{B_r \times d}$, along with $\ell_i, m_i = (0), (-\infty) \in \mathbb{R}^{B_r}$.
      \STATE Synchronize and ensure both $\mathbf{Q}_i$ and $\mathbf{K}_0$ are available in shared memory.
      \STATE With WGMMA, calculate $\mathbf{S}_{\mathrm{cur}} = \mathbf{Q}_i \mathbf{K}_0^T$. Commit this operation and await its completion.
      \STATE Release the buffer’s $0$th stage reserved for $\mathbf{K}$.
      \STATE Derive $m_{i}$, $\tilde{\mathbf{P}}_{\mathrm{cur}}$, and $\ell_{i}$ from $\mathbf{S}_{\mathrm{cur}}$ and apply scaling to $\mathbf{O}_i$.
      \FOR{$1 \le j < T_c - 1$}
        \STATE Wait until $\mathbf{K}_j$ is resident in shared memory.
        \label{code:mainloop_start}
        \STATE Perform the computation $\mathbf{S}_{\mathrm{next}} = \mathbf{Q}_i \mathbf{K}_{j}^T$ with WGMMA. Commit, but do not immediately wait. \label{code:first_wgmma}
        \STATE Pause until $\mathbf{V}_{j-1}$ is present in shared memory.
        \STATE Update $\mathbf{O}_{i} = \mathbf{O}_{i} + \tilde{\mathbf{P}}_{\mathrm{cur}} \mathbf{V}_{j-1}$ using WGMMA. Commit the computation, proceed without blocking. \label{code:second_wgmma}
        \STATE Await the completion of the WGMMA operation for $\mathbf{Q}_i \mathbf{K}_{j}^T$.
        \STATE Calculate $m_{i}$, $\tilde{\mathbf{P}}_{\mathrm{next}}$, and $\ell_{i}$ by processing $\mathbf{S}_{\mathrm{next}}$. \label{code:softmax}
        \STATE Wait for the result of WGMMA on $\tilde{\mathbf{P}}_{\mathrm{cur}} \mathbf{V}_{j-1}$, then rescale $\mathbf{O}_i$.
        \STATE Free the $(j\,\%\,s)$th stage of the $\mathbf{K}$ buffer and $(j-1\,\%\,s)$th stage of the $\mathbf{V}$ buffer, respectively.
        \STATE Assign the content of $\mathbf{S}_{\mathrm{next}}$ into $\mathbf{S}_{\mathrm{cur}}$.
        \label{code:mainloop_end}
      \ENDFOR \STATE Synchronize to ensure $\mathbf{V}_{T_c - 1}$ has been loaded into shared memory.
      \STATE Compute the final output $\mathbf{O}_{i} = \mathbf{O}_{i} + \tilde{\mathbf{P}}_{\mathrm{last}} \mathbf{V}_{T_c - 1}$ with WGMMA, committing and waiting until complete.

\STATE Epilogue: Adjust $\mathbf{O}_{i}$ according to $m_{i}$. Derive $L_{i}$ using both $m_{i}$ and $\ell_{i}$.  
Store $\mathbf{O}_{i}$ and $L_{i}$ in HBM as the $i$-th segments of $\mathbf{O}$ and $L$.

\cref{alg:flash3_wgmma} serves as a substitute for the consumer portion described in \cref{alg:flash3_wgmma_ws_only}, thereby forming the full \textsc{FlashAttention-3} procedure for FP16 calculations. Broadly, we use WGMMA to refer to asynchronous GEMM in this context. During the primary loop (spanning lines \ref{code:mainloop_start} through \ref{code:mainloop_end}), the second WGMMA call within iteration $j$ (see line \ref{code:second_wgmma}) is executed concurrently with the softmax computations belonging to the subsequent iteration, $j+1$ (line \ref{code:softmax}).

Although the pipelined approach outlined previously promises improvements in theoretical throughput, a number of practical challenges arise:

\paragraph{Compiler reordering} The pseudocode provides an idealized, sequential ordering of execution; however, in reality, the NVCC compiler frequently reschedules instructions for optimization purposes. Such instruction reordering can disrupt the planned interleaving of WGMMA and non-WGMMA operations, potentially undermining the intended pipelining benefits and even introducing anomalous behavior. Examination of the resulting SASS confirms that, in practice, the compiler still emits code with the anticipated degree of overlap (refer to Section~\ref{sec:2-stage-sass}).

\paragraph{Register pressure} Avoiding register spilling is essential for preserving high throughput. The 2-stage pipelining strategy introduces extra requirements for register storage, particularly for holding intermediate computation results and preserving inter-stage context. This is manifested most notably in the necessity to allocate space for an additional $\mathbf{S}_{\mathrm{next}}$ in registers, which incurs a per-threadblock overhead of $B_r \times B_c \times \text{sizeof}(\text{float})$. Such heightened register consumption competes directly with efforts to use larger threadblock (tile) sizes—another technique often employed for performance gains and itself demanding in terms of register resources. Consequently, practitioners must judiciously weigh these competing factors, informed by empirical profiling.

\paragraph{3-stage pipelining} Building on the aforementioned 2-stage structure, we introduce a 3-stage pipeline that further overlaps the second WGMMA invocation with the softmax computation. This refinement may yield higher utilization of Tensor Cores. Nonetheless, it imposes yet greater register requirements, as every additional pipeline stage necessitates further temporary storage. The increased complexity exacerbates the tension between maximizing tile dimensions and expanding pipeline depth. Further specifics about the 3-stage method, as well as performance evaluation, are presented in ~\cref{sec:3-stage}.

\subsection{Low-precision with FP8}
\label{sec:algofp8}

\textbf{Efficiency: layout transformations.} Executing the forward pass of \textsc{FlashAttention-3} with FP8 precision introduces unique challenges related to layout compatibility, which do not arise when using FP16.

To begin, we observe that the input tensors $\mathbf{Q}$, $\mathbf{K}$, and $\mathbf{V}$ are generally organized such that the head dimension is contiguous in memory. However, to adhere to the k-major requirement imposed by FP8 WGMMA during the second GEMM operation, it is necessary for $\mathbf{V}$—more specifically, for the tiles of $\mathbf{V}$ transferred into SMEM—to be contiguous along the sequence length dimension. Since the TMA load operation is unable to alter which dimension is contiguous, two primary strategies are available: (1) preemptively transpose $\mathbf{V}$ in GMEM, or (2) perform an in-kernel transpose on $\mathbf{V}$'s tiles once they have been loaded into SMEM. For the first approach, implementation can proceed via (1a) fusing the transpose into the epilogue of an earlier stage such as rotary embedding, or (1b) utilizing a dedicated pre-processing kernel\footnote{An optimized transpose kernel can approach device-level bandwidth in terms of speed~\citep{colfax_cutlass_transpose_2024}.} to switch the strides of the sequence length and head dimensions. Nonetheless, the fused epilogue approach (1a) poses integration challenges for standard libraries, and the standalone transpose kernel (1b) tends to be overly expensive for inference workloads, which are constrained by memory bandwidth.

For FP8 \textsc{FlashAttention-3}, we select strategy (2). Within the kernel, we exploit the LDSM (\verb|ldmatrix|) and STSM (\verb|stmatrix|) instructions to facilitate transposing. These instructions enable a warp of threads to collaboratively transfer data between SMEM and RMEM at a granularity of 128 bytes.\footnote{Although LDSM/STSM are specified in the PTX documentation as copying $8 \times 8$ matrices with 16-bit elements \cite[\S 9.7.13.4.15-16]{ptx}, it is possible to pack two 8-bit elements together, allowing LDSM/STSM to be used for FP8-level operations. However, since the transpose variants of LDSM/STSM do not support splitting packed 8-bit values, certain register-level shuffling must occur between the LDSM and STSM phases in order to execute a proper tile-wise transpose; we elide further specifics here.} These instructions are highly efficient in terms of register usage, enabling their execution within the producing warpgroup and permitting simultaneous data layout transpositions during memory copy. Furthermore, starting from the second iteration, it is possible to pipeline the transpose operation of the subsequent $\mathbf{V}$ tile so that it runs concurrently—hidden behind—the two WGMMAs operating on the previous $\mathbf{V}$ tile and the current $\mathbf{K}$ tile.

Second, it is apparent that, contrary to FP16, the FP32 accumulator utilized by an FP8 WGMMA has a memory layout that does not align with that of operand A when this operand is stored in registers. Illustrations of partial layouts for these two cases can be found in \cref{fig:rmem_accum} and \cref{fig:rmem_operand}, where register allocation per thread follows the orders indicated. By leveraging byte permute instructions, it is possible to rearrange the register contents of the first WGMMA's accumulator into a layout compatible with a subsequent WGMMA, matching also the structure of the $\mathbf{V}$ tile obtained from the in-kernel transpose. In detail, referring to \cref{fig:rmem_accum}, the transformation consists of sequencing the register contents as
$$\{ \verb|d0 d1 d4 d5 d2 d3 d6 d7| \},$$
and repeating this permutation pattern every 8 bytes. With respect to the logical organization of the $\mathbf{P}$ tile, this operation permutes its columns—so, for example, columns $0189$ become the first four columns. To ensure that WGMMA outputs the appropriate tile, the in-kernel transpose may be configured to apply a corresponding permutation to the rows of the $\mathbf{V}$ tile.\footnote{Having this extra degree of freedom when transposing in-kernel obviates the necessity for shuffle instructions to reassign register data among threads, which we previously discussed in~\citep{colfax_fp8_flashattention_2024}.}

\textbf{Accuracy: block quantization and incoherent processing.} The FP8 (e4m3) numerical format allocates just 3 bits to the mantissa and 4 bits to the exponent, leading to greater numerical inaccuracy when compared to FP16 or BF16 formats. Additionally, large-scale models often feature outlier elements~\citep{dettmers2208llm, sun2024massive} whose magnitudes greatly exceed typical values, complicating the quantization process. A common strategy is to implement per-tensor scaling~\citep{micikevicius2022fp8}, assigning a single scaling factor to each tensor (such as one each for $\mathbf{Q}$, $\mathbf{K}$, and $\mathbf{V}$). To minimize the impact of numerical errors introduced by FP8 within attention mechanisms, we utilize two specific methods:

\begin{enumerate}

\item \textbf{Block quantization}: In this approach, a single scalar value is retained for each block; thus, for the tensors $\mathbf{Q}$, $\mathbf{K}$, and $\mathbf{V}$, we partition them into blocks of size $B_r \times d$ or $B_c \times d$ and quantize each block independently. This quantization step can be merged with a preceding operation before the attention mechanism—such as rotary embedding—without incurring any runtime penalty, because rotary embedding is constrained by memory bandwidth rather than computation. Given that the \textsc{FlashAttention-3} algorithm is inherently designed to process blocks, it is possible to rescale each $\mathbf{S}$ block accordingly to compensate for the block-wise quantization, without adding computational overhead.

\item \textbf{Incoherent processing}: To mitigate the influence of outlier values, we pre-multiply both $\mathbf{Q}$ and $\mathbf{K}$ by a randomly chosen orthogonal matrix $\mathbf{M}$ before applying FP8 quantization. Because $\mathbf{M}$ satisfies $\mathbf{M} \mathbf{M}^\top = I$, the expression $(\mathbf{Q} \mathbf{M}) (\mathbf{K} \mathbf{M})^\top$ yields the same result as $\mathbf{Q} \mathbf{K}^\top$, ensuring the attention computation remains unchanged. This randomization distributes outlier components across all entries in $\mathbf{Q} \mathbf{M}$ or $\mathbf{K} \mathbf{M}$, effectively transforming each value into a weighted sum of multiple elements and thereby attenuating quantization noise. As implemented in previous works such as \citet{chee2024quip} and~\citet{tseng2024quip}, $\mathbf{M}$ is typically constructed by multiplying random diagonal matrices of $\pm 1$ with a Hadamard matrix. This structure permits matrix multiplication in $O(d \log d)$ time, as opposed to $O(d^2)$, and can also be integrated with rotary embedding without additional computational burden.
\end{enumerate}

We empirically demonstrate in \cref{sec:numerical_error} that these methods can decrease numerical error by as much as a factor of 2.6.

\section{Empirical Validation}
\label{sec:exp}

To implement \textsc{FlashAttention-3} and assess both its performance and precision, we leverage primitives provided by CUTLASS~\citep{Thakkar_CUTLASS_2023}, including the WGMMA and TMA abstractions.

\begin{itemize}

\item \textbf{Benchmarking attention.} We evaluate the execution time of \textsc{FlashAttention-3} across varying sequence lengths, making direct comparisons to the default PyTorch implementation, \textsc{FlashAttention-2}, as well as the Triton variant of \textsc{FlashAttention-2} (which leverages H100-tailored instructions), and a vendor-supplied cuDNN implementation of \textsc{FlashAttention-2} specifically optimized for H100 GPUs. Our results show that \textsc{FlashAttention-3} achieves up to 2.0$\times$ speedup over \textsc{FlashAttention-2} and is 1.5$\times$ quicker compared to the Triton-based version. On H100 hardware, \textsc{FlashAttention-3} attains up to 740 TFLOPs/s, reaching 75\% of the peak theoretical performance.
\item \textbf{Ablation study.} Our investigation demonstrates that enhancements such as warp-specialization and the GEMM-softmax pipeline are significant contributors to the overall acceleration observed with \textsc{FlashAttention-3}.
\item \textbf{Accuracy of FP8 attention.} Through experiments, we establish that the use of block quantization together with incoherent computation lowers numerical errors in FP8 \textsc{FlashAttention-3} by a factor of 2.6$\times$.
\end{itemize}

\subsection{Benchmarking Attention}
\label{subsec:benchmark_attn}

We evaluate the performance of various attention mechanisms on an H100 80GB SXM5 GPU, considering multiple configurations (with and without a causal mask, head dimensions of 64 or 128), all using FP16 precision. The benchmarking results, presented in~\cref{fig:benchmark_attn_fwd} and~\cref{fig:benchmark_attn_bwd}, indicate that \textsc{FlashAttention-3} achieves approximately 1.5-2.0$\times$ speedup over \textsc{FlashAttention-2} in the forward computation, and is about 1.5-1.75$\times$ faster in the backward computation. When contrasted with conventional attention, \textsc{FlashAttention-3} demonstrates speed improvements ranging from 3$\times$ to as high as 16$\times$. Notably, for sequence lengths of 1k tokens or greater, \textsc{FlashAttention-3} even outperforms a proprietary vendor library (cuDNN -- closed source), which has been tailored for H100 GPUs.

\paragraph{Benchmark settings:} The sequence length is adjusted across values of 512, 1k, up to 16k, while the batch size is chosen such that the aggregate token count remains fixed at 16k. The hidden dimension is fixed at 2048, and the head dimension is varied among 64, 128, and 256 (corresponding to 32, 16, or 8 attention heads, respectively). FLOPs for the forward pass are determined using:

\begin{equation*}
  4 \cdot \text{seqlen}^2 \cdot \text{head dimension} \cdot \text{number of heads}.
\end{equation*}

When employing causal masking, we halve this quantity to reflect that roughly half of the elements are actually computed. For estimating the FLOPs required during the backward pass, we take the FLOPs from the forward pass and multiply by 2.5, given that the backward pass involves five matrix multiplications (due to recomputation), compared to two in the forward pass.

Additionally, we evaluate the forward pass runtime using FP8 under comparable conditions. The outcomes for headdim 256 are presented in ~\cref{fig:benchmark_attn_fp8}, with comprehensive results detailed in \cref{sec:benchmark_attn_fp8_full}.

\subsection{Ablation Study: 2-Stage Pipelining Experiments}
\label{sec:2-stage-pipelining-experiments}

We conduct an ablation study on both the 2-stage WGMMA-softmax pipelining and warp-specialization techniques for non-causal FP16 \textsc{FlashAttention-3}, maintaining constant parameters $\{\text{batch}, \text{seqlen}, \text{nheads}, \text{hdim}\} = \{ 4, 8448, 16, 128\}$. As shown in~\cref{table:ablation_pipelining}, our algorithmic strategies—specifically, incorporating asynchrony via warp-specialization and allowing GEMM and softmax operations to overlap—produce notable performance gains, boosting throughput from 570 to 661 TFLOPs.

\subsection{Numerical Error Validation}
\label{sec:numerical_error}

Given the ongoing discussion around the numerical error~\citep{golden2024flash} associated with \textsc{FlashAttention}, we evaluate the numerical accuracy of \textsc{FlashAttention-2}, \textsc{FlashAttention-3}, and a conventional attention implementation, using an FP64-based reference as a benchmark. To better reflect the presence of extreme features and activations observed in LLMs~\citep{dettmers2208llm, sun2024massive}, we sample entries of $\mathbf{Q}, \mathbf{K}, \mathbf{V}$ from the following distribution:

\begin{equation*}
  \mathcal{N}(0, 1) + \mathcal{N}(0, 100) \cdot \mathrm{Bernoulli}(0.001).
\end{equation*}

In this setup, every element is sampled from a normal distribution with mean zero and unit standard deviation, except for 0.1\% of elements, to which we add an independent Gaussian noise component with a standard deviation of 10. The root mean squared error (RMSE) corresponding to each method is reported in~\cref{table:numerical_error}. Using FP16 precision, both \textsc{FlashAttention-2} and \textsc{FlashAttention-3} yield RMSE values that are 1.7$\times$ smaller than those from the conventional approach, since intermediate computations, specifically the softmax outputs, are retained in FP32. The baseline attention algorithm in FP8 leverages per-tensor scaling, accumulates matmul operations in FP32, and holds softmax intermediates in FP16. Owing to block quantization and incoherent computation techniques, \textsc{FlashAttention-3} operating in FP8 achieves 2.6$\times$ the accuracy of this baseline.

\section{Dicussion, Limitations, Conclusion}
\label{sec:discussion}

Through the development of \textsc{FlashAttention-3}, we illustrate that leveraging novel programming strategies alongside hardware advancements like asynchrony and low-precision arithmetic can significantly enhance both the speed and precision of attention mechanisms. Our method achieves a 1.5-2.0$\times$ speedup relative to \textsc{FlashAttention-2}, while also decreasing FP8 numerical error by a factor of 2.6$\times$ in comparison with traditional per-tensor quantization approaches. There remain some challenges, which we intend to address in future research: further tuning for large language model (LLM) inference, introducing a persistent kernel framework within the FP8 kernel,\footnote{In our experiments, the FP16 \textsc{FlashAttention-3} version incorporates a persistent kernel and dynamic load balancing, whereas FP8 \textsc{FlashAttention-3} does not. This partially accounts for the lower performance of FP8 \textsc{FlashAttention-3} in scenarios with short sequence lengths and causal masking, relative to FP8 cuDNN kernels.} and thoroughly investigating the ramifications of low-precision attention for extensive model training. Although the current study emphasizes Hopper GPUs, we anticipate that the techniques presented are adaptable to a range of other hardware accelerators. Ultimately, we hope that the development of more efficient and accurate attention primitives will pave the way for progress in long-context processing tasks.

\end{document}